% =====================================================================
% 退避した内容（japanese/main_JP.tex，日本語版）
% 2026-10-29 の ic-slim 再構成（ブランチ mcfg-ic-slim-20261029）で本文から外したもの．
% 元の版：コミット 66897a1 の main_JP.tex（全文はブランチ archive/mcfg-pre-slim-full-20261029 にも保存）．
% 各ブロックは元の文面のままの複製．英語版の対応ブロックは parked_EN.tex の同じタグ [P*]．
% 本文に \input しないこと（ラベルが衝突する）．状態・依存関係・復元方法は README.md を参照．
% =====================================================================


% ---------------------------------------------------------------------
% [P1] 付録H（補足的な結果）全体
% ---------------------------------------------------------------------
\section{補足的な結果}
\label{app:supplementary}

この付録には，本論を補うが本論では用いない結果をまとめる．すなわち，短い証拠集合の効率的な計算，分布的代表からの特徴標本（アーベル型フィルタ付き標的を含む），\ref{subsec:baseline}節の素性族に対する明示的なデータ予算，群値の一致の言語族，本質的にランク2を要する標的である．記号は\ref{sec:char-data}節に従う．特に$C_s(G)=|V||\mathsf S^\sharp|^f+|P||\mathsf S^\sharp|^{fr_{\max}}\le|\mathsf S^\sharp|^{fr_{\max}}\|G\|$である．

\subsection{短い特徴データの計算}
\label{subsec:effective-char-data}

$G$をファンアウト$\le f$，ランク$\le r_{\max}$の縮約済み良形文法とする．型付き規則$R:X\to\rho(Y_1,\ldots,Y_q)$について，$\rho$が書く終端記号の数を$\mathrm{ter}_R$，$\operatorname{lhs}R:=X$とし，$\widetilde G_0$の型付き規則（ランク0を含む）の集合を$\mathcal R$とする．
\[
 \ell_{\mathrm{in}}(X):=\min_{\tuple u\in L_X}\|\tuple u\|_1,\qquad
 \ell_{\mathrm{out}}(X):=\min\{|\chi|_{\mathrm{fix}}:\chi\text{は}X\text{について導出的}\}
\]
とおく（導出的は補題~\ref{lem:thickness-equivalence}(iv)の意味であり，$\widetilde S\to X$が規則なら$\ell_{\mathrm{out}}(X)=0$）．さらに
\[
 b(R):=\mathrm{ter}_R+\sum_{j=1}^q\ell_{\mathrm{in}}(Y_j),\qquad
 \mathsf B_{\mathrm{rule}}(G,s):=
 \sum_{R\in\mathcal R}\bigl(1+\ell_{\mathrm{out}}(\operatorname{lhs}R)+b(R)\bigr)
\]
とおく．

\begin{proposition}[短い特徴データの計算と規則のみのデータ予算]
\label{prop:effective-char-data}
\label{prop:computed-witness-budget}
\label{prop:rule-only-witnesses}
$f$，$r_{\max}$，$s$を固定する．$G$が与えられたとき，次は$\|G\|$の多項式時間で計算できる（数は二進表記）．
\begin{enumerate}[label=\textup{(\roman*)},leftmargin=*]
\item $\widetilde G_0$，全ての値$\ell_{\mathrm{in}}(X)$，$\ell_{\mathrm{out}}(X)$，$\mathsf B_{\mathrm{rule}}(G,s)$，および$G$の各規則$R$の$\theta_G(R)$（よって$\theta_G$と$t_G$）．
\item 各型付き非終端記号$X$について，再帰$\omega(X):=\rho_X(\omega(Y_1),\ldots,\omega(Y_q))$が整礎であり，$\|\omega(X)\|_1=\ell_{\mathrm{in}}(X)$なる$\omega(X)\in L_X$を与えるような規則$R_X:X\to\rho_X(Y_1,\ldots,Y_q)$，および$|\chi(X)|_{\mathrm{fix}}=\ell_{\mathrm{out}}(X)$なる導出的文脈$\chi(X)$（$\widetilde S\to X$が規則なら$\blank_1$）．いずれも$\|G\|$の多項式サイズの導出DAGで与えられる．
\end{enumerate}
これらの選択について
\[
 \begin{gathered}
 \Wit(\widetilde G_0)=\bigl\{\chi(\operatorname{lhs}R)[\rho_R(\omega(Y_1),\ldots,\omega(Y_q))]:
 R\in\mathcal R\bigr\},\\
 |\Wit(\widetilde G_0)|\le|\mathcal R|,\qquad
 \|\Wit(\widetilde G_0)\|_+\le\mathsf B_{\mathrm{rule}}(G,s)
 \end{gathered}
\]
であり，語は$\|G\|+\|\Wit(\widetilde G_0)\|_+$の多項式時間で書き出せる．\ref{subsec:thickness-contexts}--\ref{subsec:algebraic-typed-data}節，\ref{subsec:rank-one}節，付録~\ref{app:local-group}で$\Wit(\widetilde G_0)$について示した全ての上界はこの集合について成り立つ．特に，$L\in\Cmcf{f}{s}$について$L(G)=L\cap\Sigma^+$なら，これは特徴標本である．
\end{proposition}

\begin{proof}
$f,r_{\max},s$を固定すれば，完全精密化の記号と規則は$O(C_s(G))$個であり，刈り込みは線形時間の生産性・到達可能性の計算である．

\emph{内側の値．}導出に重み（終端長，節点数）を辞書式順序で与える．規則$R$は子の重み$(x_j,n_j)$を$(\mathrm{ter}_R,1)+\sum_j(x_j,n_j)$に写し，これは各引数より真に大きい．よってNederhofの重み付き演繹の形~\cite{Nederhof2003}によるKnuthのDijkstra法の一般化（\cite{RamalingamReps1996}も参照）は，各型付き非終端記号の最小重みと，子が$X$より先に確定する最小化規則$R_X$を計算する．この順序により$\omega$の再帰は整礎であり，選んだ規則はDAGをなし，重みの第一成分は$\ell_{\mathrm{in}}(X)=\|\omega(X)\|_1$である．

\emph{外側の値．}各型付き規則$R:X\to\rho(Y_1,\ldots,Y_q)$について重み$\mathrm{ter}_R+\sum_{i\ne j}\ell_{\mathrm{in}}(Y_i)$の辺$X\to Y_j$をもつグラフ上で，$\widetilde S\to X$なる$X$を距離0の始点としてDijkstra法を実行する．始点から$X$への経路は，テンプレートを合成し兄弟を$\omega(Y_i)$で埋めると，その固定長の$X$の文脈を与える．これは補題~\ref{lem:filling-identity}により許容的であり，成功導出から$X$-部分木を削除して得られるので導出的である．逆に$X$の導出的文脈は根から$X$への経路から生じ，その兄弟の生成語は$\ell_{\mathrm{in}}$以上の長さをもつ．よって距離は$\ell_{\mathrm{out}}(X)$であり，最短経路木と内側のDAGが$\chi(X)$を表す．同じ二つの計算を$G$に適用すれば，導出の外側と内側は独立に選べるので，$t_G=\max_A\ell_{\mathrm{in}}(A)$および$R:A\to\rho(B_1,\ldots,B_q)$について$\theta_G(R)=\ell_{\mathrm{out}}(A)+\mathrm{ter}_R+\sum_j\ell_{\mathrm{in}}(B_j)$が得られる．全ての値は$\|G\|$について高々指数的なので，二進長は多項式である．

\emph{証拠集合と予算．}$\omega$の定義により$\chi(X)[\omega(X)]$は$R_X$の規則語であり，定義~\ref{def:witness-set}の他の語も全て規則語である．$R$の規則語の長さはちょうど$\ell_{\mathrm{out}}(\operatorname{lhs}R)+b(R)$であり，$\mathcal R$にわたって$|w|+1$を足すと$\mathsf B_{\mathrm{rule}}$になる．相異なる語に移れば値は下がるだけである．DAGの展開には表す語の長さに線形な時間がかかる．

\emph{上界の継承．}引用した各上界は，\emph{ある}$\omega(X)\in L_X$と導出的文脈$E_B$（補題~\ref{lem:thickness-equivalence}(iv)）を示すことで証明されている．計算した最小値はそれより長くなく，語の数も変わらない．特徴標本であることは補題~\ref{lem:witness-length-template}(iv)による．
\end{proof}

予算$\mathsf B_{\mathrm{rule}}$は表示に依存する\emph{上界の証明書}であり，最悪の場合の上界よりはるかに小さいことが多いが，最小の特徴標本のサイズではない．\ref{subsec:rank-one}節の$\{\mathtt c^{2^n}\}$では語は$\|G\|$より指数的に長いが，そのDAGは多項式サイズである．

\subsection{分布的代表とアーベル型フィルタ付き標的}
\label{subsec:short-distributional-representatives}

命題~\ref{prop:completeness}は，代表が自身の非終端記号から導出できることを要求しない．したがって型付き分布同値類全体の最短元を用いてよい．$G$を$L\in\Cmcf{f}{s}$なる$L\cap\Sigma^+$のファンアウト$\le f$，ランク$\le r$の縮約済み良形表示とする．型付き非終端記号$X$について，$L_X$を含む$\equiv_{L,s}$-同値類を$\operatorname{Cl}_{L,s}(X)$とし（補題~\ref{lem:typed-homogeneity}），
\[
 \kappa_G(X):=\min_{\tuple v\in\operatorname{Cl}_{L,s}(X)}\|\tuple v\|_1\le\ell_{\mathrm{in}}(X),\qquad
 \kappa_G:=\max_X\kappa_G(X)
\]
とおく．

\begin{proposition}[緩和した分布的証拠]
\label{prop:mcfg-relaxed-witnesses}
各型付き非終端記号$X=B_{\mathbf t}$について$\|\iota(X)\|_1=\kappa_G(X)$なる$\iota(X)\in\operatorname{Cl}_{L,s}(X)$を選び，$\chi(X):=E_B$（補題~\ref{lem:thickness-equivalence}(iv)）とする．全ての型付き非終端記号と型付き規則$X\to\rho(Y_1,\ldots,Y_q)$にわたる語$\chi(X)[\iota(X)]$と$\chi(X)[\rho(\iota(Y_1),\ldots,\iota(Y_q))]$は集合$K^\dagger\subseteq L$をなし，$K^\dagger\cup(\{\eps\}\cap L)$は$\mathcal B^{\le r}_{f,s}$と$\mathcal B_{f,s}$の特徴標本であり，
\[
 |K^\dagger|\le C_s(G),\qquad
 \|K^\dagger\|_+\le C_s(G)\bigl(1+2\bar\theta_G+r\kappa_G\bigr)
\]
である．
\end{proposition}

\begin{proof}
各$X$について$\tuple u_X\in L_X$を固定する．$\iota(X)\equiv_{L,s}\tuple u_X$であり$E_B$は$\tuple u_X$を受理するので，$\iota(X)$も受理する．よって$\chi(X)[\iota(X)]\in L$である．補題~\ref{lem:witnessed-composition}により$\rho(\iota(Y_1),\ldots,\iota(Y_q))\equiv_{L,s}\rho(\tuple u_{Y_1},\ldots,\tuple u_{Y_q})\in L_X$なので，規則語も$L$に属する．各$\iota(X)$は$\widehat s$-ベクトル$\mathbf t$と非空成分をもち，$\chi(S_{\mathbf t})=E_S=\blank_1$なので，この$\iota$について命題~\ref{prop:completeness}が適用でき，$K\supseteq K^\dagger$について$L\cap\Sigma^+\subseteq L(\mathcal B(K))$を得る．$K\subseteq L$なら逆向きは命題~\ref{prop:reconstruction-soundness}により，$\eps$はStart規則が扱う．長さの上界は$b=\kappa_G$とした補題~\ref{lem:witness-length-template}(ii)であり，その証明は$\|\iota(X)\|_1\le b$しか用いない．
\end{proof}

命題~\ref{prop:effective-char-data}と異なり，語$\chi(X)[\iota(X)]$は一般には省けない．$f,r,s$を固定したとき，$\|G_n\|\bar\theta_{G_n}=n^{O(1)}$でありながら超多項式の特徴データを要する族$G_n$は，$\kappa_{G_n}=n^{\omega(1)}$でなければならない．これは必要条件であり，そのような族が存在することの証明ではない．

\begin{theorem}[アーベル型フィルタ付き標的の多項式データ]
\label{thm:mcfg-abelian-filter-data}
空でないアルファベット$\Sigma$，型付け$s$，$f,r\ge1$を固定する．$\Sigma,s,f,r$のみに依存する定数$C_{\mathrm{ab}}$が存在し，任意のアーベル群$\mathsf A$，準同型$\varphi:\Sigma^*\to\mathsf A$，部分群$\mathsf H\le\mathsf A$，$\gamma\in\mathsf A$，$\mathsf{Acc}\subseteq\mathsf S^\sharp$について次が成り立つ．$L:=\widehat s^{-1}(\mathsf{Acc})\cap\varphi^{-1}(\gamma+\mathsf H)$とし，$L\cap\Sigma^+$がファンアウト$\le f$，ランク$\le r$の縮約済み良形表示$G$をもつなら，$L\in\Cmcf{f}{s}$であり，$\mathcal B^{\le r}_{f,s}$と$\mathcal B_{f,s}$は
\[
 \|K\|_+\le C_{\mathrm{ab}}\,\bigl(\|G\|\bar\theta_G\bigr)^{3|\Sigma|+2}
\]
なる共通の特徴標本$K$をもつ．
\end{theorem}

\begin{proof}
命題~\ref{prop:abelian-coset-filter}により$L\in\Cmcf{f}{s}$である．命題~\ref{prop:mcfg-relaxed-witnesses}と$C_s(G)\le|\mathsf S^\sharp|^{fr}\|G\|$により，$\Theta:=(r+1)t_G+\|G\|\le(r+2)\|G\|\bar\theta_G$と$\Sigma,s,f$のみに依存する$C'$について$\kappa_G\le C'(\Theta+1)^{3|\Sigma|+1}$を示せば十分である．

\emph{許容されるParikhシフトの格子．}$\psi:\Sigma^*\to\mathbb N^\Sigma$をParikh写像$\psi(w)=(|w|_a)_{a\in\Sigma}$とし，$\Psi(\tuple u):=\sum_i\psi(u_i)$とおき，$\varphi=\overline\varphi\circ\psi$なる準同型を$\overline\varphi:\mathbb Z^\Sigma\to\mathsf A$，$\Upsilon:=\overline\varphi^{-1}(\mathsf H)$とする．各$C\in V$について$\|\omega_C\|_1\le t_G$なる$\omega_C\in L_C(G)$を固定する．$\tuple u\in L_C(G)$について，文脈$E_C$は$\tuple u$と$\omega_C$を受理し，$L\subseteq\varphi^{-1}(\gamma+\mathsf H)$なので
\[
 \Psi(\tuple u)-\Psi(\omega_C)\in\Upsilon\tag{$*$}
\]
である．規則$R:C\to\rho(B_1,\ldots,B_q)$について，終端文字のParikhベクトルを$\psi_R\in\mathbb N^\Sigma$とする．$\delta_R:=\Psi(\omega_C)-\psi_R-\sum_j\Psi(\omega_{B_j})$とおく．$(*)$を$\rho(\omega_{B_1},\ldots,\omega_{B_q})$に適用すれば$\delta_R\in\Upsilon$であり，$\|\delta_R\|_1\le \Theta$である．$\Gamma_\Theta:=\{z\in\Upsilon:\|z\|_1\le \Theta\}$，$\Upsilon_\Theta:=\langle\Gamma_\Theta\rangle$とする．$\tuple u_j\in L_{B_j}(G)$による$\tuple u=\rho(\tuple u_1,\ldots,\tuple u_q)$について$\Psi(\tuple u)-\Psi(\omega_C)=\sum_j(\Psi(\tuple u_j)-\Psi(\omega_{B_j}))-\delta_R$なので，導出に関する帰納法により，全ての$\tuple u\in L_C(G)$について$\Psi(\tuple u)-\Psi(\omega_C)\in\Upsilon_\Theta$である．

\emph{短い代表．}$X=C_{\mathbf t}$をアリティ$d\le f$，$\tuple u\in L_X$とする．各$i$について，正則言語$\widehat s^{-1}(t_i)\cap\Sigma^+$のParikh像（Parikhの定理~\cite{Parikh1966}により半線形である）の半線形表現（有限個の線形集合の合併）を一度だけ固定し（$G$は$\lambda$-自由なので$L_X$のタプルの各成分は空でなく，その型は$\mathbf 1_\sharp$にならないから$t_i\in\mathsf S$である），$\psi(u_i)$を含むその線形集合$\mathbf o_i+\sum_{p\in P_i}\mathbb Np$を選ぶ．生成元の数と大きさは$\Sigma$と$s$のみに依存する．$m_z\in\mathbb Z$を用いて$\Psi(\tuple u)-\Psi(\omega_C)=\sum_{z\in\Gamma_\Theta}m_zz$と書く．このとき連立方程式
\[
 \sum_{i=1}^d\sum_{p\in P_i}\eta_{i,p}\,p-\sum_{z\in\Gamma_\Theta}(\mu^+_z-\mu^-_z)\,z
 =\Psi(\omega_C)-\sum_{i=1}^d\mathbf o_i
\]
は非負整数解をもつ（$\psi(u_i)$の係数と$\mu^\pm_z:=\max\{\pm m_z,0\}$）．方程式は$\bar m=|\Sigma|$本，未知数は$\bar n\le D+2(2\Theta+1)^{|\Sigma|}$個，係数の絶対値は$a\le D(\Theta+1)$以下であり，$D$は$\Sigma,s,f$のみに依存する．Papadimitriouの上界~\cite{Papadimitriou1981}により，全成分が$\bar n(\bar m a)^{2\bar m+1}=O\bigl((\Theta+1)^{3|\Sigma|+1}\bigr)$以下の非負整数解が存在する．各$i$についてベクトル$\mathbf o_i+\sum_p\eta_{i,p}p$はある空でない$v_i\in\widehat s^{-1}(t_i)$のParikhベクトルであり，$\tuple v:=(v_1,\ldots,v_d)$とおく．すると$\|\tuple v\|_1=O((\Theta+1)^{3|\Sigma|+1})$，$\widehat s^{(d)}(\tuple v)=\mathbf t$，$\Psi(\tuple v)-\Psi(\tuple u)\in\Upsilon_\Theta\subseteq\Upsilon$である．任意の許容文脈$E$について，$E[\tuple v]$と$E[\tuple u]$は同じ$\widehat s$-値をもち，$\varphi(E[\tuple v])-\varphi(E[\tuple u])=\overline\varphi(\Psi(\tuple v)-\Psi(\tuple u))\in\mathsf H$である．よって$E[\tuple v]\in L$と$E[\tuple u]\in L$は同値であり，$\tuple v\equiv_{L,s}\tuple u$かつ$\kappa_G(X)\le\|\tuple v\|_1$である．
\end{proof}

この定理は$\mathsf S\notin\mathbf{LG}$でも成り立つ．仮定は型付けではなく標的に関するものだからである．実効的に与えられない群$\mathsf A$について多項式時間の構成は主張しない．

\subsection{素性族の明示的なデータ予算}
\label{app:feature-budget}

記号は\ref{subsec:baseline}節に従う．

\begin{proposition}[共通の素性型付けに対する明示的なデータ予算]
\label{prop:feature-family-witness-budget}
$f\ge2$，$R\ne\emptyset$とする．作用素$\mathcal B^{\le2}_{f,s}$は$\|K_R\|_+=36|R|+22$なる特徴標本$K_R\subseteq L_R$をもつ（$R=\emptyset$なら空標本が特徴標本）．逆に，集合駆動型作用素$\mathcal D$が\emph{全ての}$L_{R'}$（$R'\subseteq\{0,1\}^2$）について特徴標本をもつなら，$\mathcal D$に関する$L_R$の任意の特徴標本$K$は，各$(i,j)\in R$について標識$\mathtt p_i,\mathtt q_j$をもつ語を含む．よって$\|K\|_+\ge7|R|$である．
\end{proposition}

\begin{proof}
長さ$2m+2n+2$の語$w^{m,n}_{ij}:=\mathtt a^m\mathtt p_i\mathtt b^n\mathtt c^m\mathtt q_j\mathtt d^n$を用い，$(i_0,j_0)\in R$を固定して
\[
 K_R:=\{w^{m,n}_{ij}:(i,j)\in R,\ m,n\in\{1,2\}\}\cup
 \{w^{3,1}_{i_0j_0},w^{1,3}_{i_0j_0}\}
\]
とおく．符号化コストは一対あたり$7+9+9+11=36$，追加の二語で$11+11$である．

\emph{型付き精密化．}$R$で添字付けられた開始規則をもつ命題~\ref{prop:feature-cross-typed}の文法$G$について，$(\mathtt a^m,\mathtt c^m)$の$s$-型は$m=1$かどうかだけで決まる（短い語は正確に記録される）．よって$A$は二つの型付きコピー$A_1,A_+$をもち，$B$も同様に$B_1,B_+$をもつ．$w^{m,n}_{ij}$の型は$(i,j)$で決まるので，各$(i,j)\in R$について型付き開始非終端記号$S_{ij}$が一つある．$\rho_A:=(\mathtt a\xi_1^1,\mathtt c\xi_1^2)$とすると，型付き規則は$A_1\to(\mathtt a,\mathtt c)$，$A_+\to\rho_A(A_1)$，$A_+\to\rho_A(A_+)$，その$B$版，および各$(i,j)\in R$についての四つの開始規則$S_{ij}\to\rho_{ij}(A_\bullet,B_\circ)$である．

\emph{証拠．}$\iota(A_1)=(\mathtt a,\mathtt c)$，$\iota(A_+)=(\mathtt a^2,\mathtt c^2)$，$\iota(B_1)=(\mathtt b,\mathtt d)$，$\iota(B_+)=(\mathtt b^2,\mathtt d^2)$，$\iota(S_{ij})=w^{1,1}_{ij}$，文脈$\chi(A_\bullet)=\blank_1\mathtt p_{i_0}\mathtt b\blank_2\mathtt q_{j_0}\mathtt d$，$\chi(B_\circ)=\mathtt a\mathtt p_{i_0}\blank_1\mathtt c\mathtt q_{j_0}\blank_2$，$\chi(S_{ij})=\blank_1$とする．命題~\ref{prop:completeness}が要求する語は，$A_1$，$B_1$とその定数規則には$w^{1,1}_{i_0j_0}$，$A_+$，$B_+$と$A_1$，$B_1$からの規則には$w^{2,1}_{i_0j_0}$と$w^{1,2}_{i_0j_0}$，$A_+$，$B_+$からの規則には$w^{3,1}_{i_0j_0}$と$w^{1,3}_{i_0j_0}$，$S_{ij}$とその四つの規則には$w^{m,n}_{ij}$（$m,n\in\{1,2\}$）である．いずれも$K_R$に属する．$K_R\subseteq L_R\in\Cmcf{f}{s}$なので，命題~\ref{prop:completeness}と命題~\ref{prop:reconstruction-soundness}により$K_R$は$\mathcal B^{\le2}_{f,s}$の特徴標本である．

\emph{下界．}$K$が標識$\mathtt p_i,\mathtt q_j$をもつ語を含まなければ，$R'=R\setminus\{(i,j)\}$について$K\subseteq L_{R'}$である．$L_{R'}$の特徴標本$F$について，集合$K\cup F$は$K$と$L_R$の間にあり，かつ$F$と$L_{R'}$の間にあるので，$L(\mathcal D(K\cup F))$は$L_R\ne L_{R'}$の両方に等しくなってしまう．相異なる対には相異なる語が必要であり，各語の長さは6以上である．
\end{proof}

この例は有限の格認可や素性の適合性の制約を模したものであり，特定の自然言語についての主張はしない．関係$R$は標的ごとに変わるが，準同型$g$と作用素は変わらない．

\subsection{群値の一致}
\label{app:group-agreement}

${\mathbb G}$を有限群，$\Delta_1$をその文字像が準同型$g_1:\Delta_1^*\to {\mathbb G}$の下で${\mathbb G}$を生成する有限アルファベットとし，$\Delta_2=\{\bar\delta:\delta\in\Delta_1\}$を互いに素なコピーとする．
$\Sigma_{\mathrm{agr}}:=\{\mathtt a,\mathtt b,\mathtt c,\mathtt d\}\uplus\Delta_1\uplus\Delta_2$とおき，準同型$g:\Sigma_{\mathrm{agr}}^*\to {\mathbb G}\times {\mathbb G}$を
\[
 g(\delta)=(g_1(\delta),1),\qquad
 g(\bar\delta)=(1,g_1(\delta)),\qquad
 g(\mathtt a)=g(\mathtt b)=g(\mathtt c)=g(\mathtt d)=(1,1)
\]
で定める．$R\subseteq {\mathbb G}\times {\mathbb G}$について
\[
 L_R^{\mathbb G}:=\{\mathtt a^m u\mathtt b^n\mathtt c^m v\mathtt d^n:
 m,n\ge1,\ u\in\Delta_1^+,\ v\in\Delta_2^+,\ g(uv)\in R\}
\]
とし，対角関係に対する言語を
$L_{\mathrm{agr}}^{\mathbb G}:=L_{\{(x,x):x\in {\mathbb G}\}}^{\mathbb G}$と書く．像への値域制限を行った正語型付け
$s_{\mathrm{agr}}:=(s_{1,1},g)|_{\Sigma_{\mathrm{agr}}^+}$
を固定する．この型付けの群成分は${\mathbb G}$ではなく${\mathbb G}\times {\mathbb G}$である．

\begin{proposition}[順序付き群値をもつ素性の一致]
\label{prop:group-agreement}
\begin{enumerate}[label=\textup{(\roman*)},leftmargin=*]
\item 任意の$R\subseteq {\mathbb G}\times {\mathbb G}$について，全ての$f\ge2$で
$L_R^{\mathbb G}\in\mathsf{MCFL}(2,2)\cap\Cmcf{f}{s_{\mathrm{agr}}}$である．$R\ne\emptyset$なら$L_R^{\mathbb G}$は文脈自由でない．
\item 交換子部分群$[{\mathbb G},{\mathbb G}]$が中心$Z({\mathbb G})$に含まれないなら，任意の$f\ge2$および任意の型付け
$s'\preceq(s_{k,\ell},g')|_{\Sigma_{\mathrm{agr}}^+}$について
$L_{\mathrm{agr}}^{\mathbb G}\notin\Cmcf{f}{s'}$である．ここで$k,\ell\ge0$で，$g'$は有限\emph{アーベル}群への準同型とする．特に$L_{\mathrm{agr}}^{\mathbb G}\notin\YSub{f}$である．
\end{enumerate}
${\mathbb G}=S_3$では\textup{(ii)}の仮定が成り立つ．一方，定理~\ref{thm:window-group-polydata}は窓・$(S_3\times S_3)$型付けのクラス$\Cmcf{f}{s_{\mathrm{agr}}}$全体に一様に適用できる．
\end{proposition}

\begin{proof}
(i) 各$x\in {\mathbb G}$について，正則言語$g_1^{-1}(x)\cap\Delta_1^+$の良形な右線形文法$U_x$をとり，その$\Delta_2$上のコピーを生成する文法を$\bar U_x$とする．全ての群元は非空正語で表現できるので，これらは存在する．ファンアウト2の規則
\[
\begin{aligned}
 P&\to(\xi_1^1,\xi_2^1)(U_x,\bar U_y) &&((x,y)\in R),\\
 C&\to(\mathtt a,\mathtt c)
      \mid(\mathtt a\xi_1^1,\mathtt c\xi_1^2)(C),\\
 D&\to(\mathtt b,\mathtt d)
      \mid(\mathtt b\xi_1^1,\mathtt d\xi_1^2)(D),\\
 Y&\to(\xi_1^1\xi_2^1,\xi_1^2\xi_2^2)(C,P),\\
 S&\to(\xi_1^1\xi_2^1\xi_1^2\xi_2^2)(Y,D)
\end{aligned}
\]
によって$L_R^{\mathbb G}$を生成でき，ファンアウトと規則ランクはいずれも高々2である．$Y$と$S$の各成分で隣接する変数は異なる子に属するので，文法は良形である．$R=\emptyset$なら言語は空である．そうでなければ$(x,y)\in R$をとり，$g_1(p)=x$，$g_1(q)=y$なる非空語$p,q\in\Delta_1^+$を選ぶ．正則言語
$\mathtt a^+p\mathtt b^+\mathtt c^+\bar q\mathtt d^+$
との共通部分をとり，素性ブロックを消去すると
$\{\mathtt a^m\mathtt b^n\mathtt c^m\mathtt d^n:m,n\ge1\}$
が得られるので，$L_R^{\mathbb G}$は文脈自由でない．

型付き代入可能性について，
$L_R^{\mathbb G}=Q\cap g^{-1}(R)\cap\varphi^{-1}(0)$と書く．ここで
$Q=\mathtt a^+\Delta_1^+\mathtt b^+\mathtt c^+\Delta_2^+\mathtt d^+$
は狭義2-局所であり，$\varphi:\Sigma_{\mathrm{agr}}^*\to\mathbb Z^2$は
$\mathtt a,\mathtt c,\mathtt b,\mathtt d$を順に
$(1,0),(-1,0),(0,1),(0,-1)$に写し，全ての素性文字を$0$に写す．$Q$には補題~\ref{lem:sl2-boundary-control}が適用でき，$g^{-1}(R)$は$s_{\mathrm{agr}}$-同値類の和集合であり，アーベルな条件と共通部分は命題~\ref{prop:abelian-coset-filter}で扱える．よって$L_R^{\mathbb G}\in\Cmcf{f}{s_{\mathrm{agr}}}$である．

(ii) $[{\mathbb G},{\mathbb G}]\not\subseteq Z({\mathbb G})$なので，中心的でない交換子
$z:=yx y^{-1}x^{-1}$が存在する．$hz\ne zh$なる$h\in {\mathbb G}$を選び，$g_1$-値がそれぞれ$x,y,h$である非空語$p,q,w\in\Delta_1^+$をとる．$\hat k\ge\max\{k+\ell,1\}$，$\delta\in\Delta_1$を選んで$\mathit{pad}:=\delta^{\hat k|{\mathbb G}|}$とおく．
\[
 u=\mathit{pad} p q q^{|{\mathbb G}|-1}p^{|{\mathbb G}|-1}\mathit{pad},\qquad
 u'=\mathit{pad} q p q^{|{\mathbb G}|-1}p^{|{\mathbb G}|-1}\mathit{pad}
\]
とする．このとき$g_1(u)=1$，$g_1(u')=z$であるが，$u,u'$のParikhベクトルは一致し，長さ$k$の接頭辞と長さ$\ell$の接尾辞も一致する．したがって，命題の任意の窓・アーベル型付けは$s'(u)=s'(u')$を満たし，コピーについても$s'(\bar u)=s'(\bar u')$である．タプル$\mathbf u:=(u,\bar u)$と$\mathbf u':=(u',\bar u')$の$s'$-型は等しい．文脈
$E:=\mathtt a\blank_1\mathtt{bc}\blank_2\mathtt d$
は，それぞれの群値が$(1,1)$と$(z,z)$なので両方を受理する．さらに
$F:=\mathtt a w\blank_1\mathtt{bc}\blank_2\bar w\mathtt d$
とおくと，$F[\mathbf u]\in L_{\mathrm{agr}}^{\mathbb G}$の群値は$(h,h)$だが，
$F[\mathbf u']\notin L_{\mathrm{agr}}^{\mathbb G}$の群値は$(hz,zh)$であり，$hz\ne zh$である．よって共通文脈$E$が存在しても分布は一致せず，アリティ2での$(f,s')$-代入可能性に反する．自明型付けも特別な場合として含まれる．${\mathbb G}=S_3$では$[S_3,S_3]=A_3$は中心に含まれない．
\end{proof}

\begin{remark}
上の分離文脈は中心的でない交換子を用いている．$Q_8$のように交換子部分群が中心に含まれる群について，同じ主張が成立するとはここでは主張しない．
\end{remark}

\subsection{本質的にランク2を要する境界型付き標的}
\label{subsec:intrinsic-ranktwo}

$\Sigma_{\mathrm{intr}}:=\{\mathtt a,\mathtt b,\mathtt l,\mathtt r,\mathtt p,\mathtt q,\mathtt c,\mathtt d,\mathtt\#,\mathtt\$,\mathtt u,\mathtt v,\mathtt w\}$，$P_{\mathrm{intr}}:=\{\mathtt{lpc},\mathtt{lpd},\mathtt{lqc},\mathtt{rpc},\mathtt{rpd}\}$，
\[
 \mathrm{MIX}_2^+:=\{x\in\{\mathtt a,\mathtt b\}^+:|x|_{\mathtt a}=|x|_{\mathtt b}\},\quad
 T_3^+:=\{\mathtt u^n\mathtt v^n\mathtt w^n:n\ge1\},\quad
 L_{\mathrm{intr}}:=P_{\mathrm{intr}}\,\mathtt\#\,\mathrm{MIX}_2^+\,\mathtt\$\,T_3^+
\]
とする．因子$\mathrm{MIX}_2^+$は規則ランク2を，$T_3^+$はファンアウト2を強制し，接頭部$P_{\mathrm{intr}}$は標的を$\YSub{f}$から分離する．

EDT0L系は，終端アルファベットを含む有限の作業アルファベット，公理，および各記号に一つの語を割り当てる置換である有限個の\emph{表}からなる．公理に表の有限列を適用して得られる終端語を生成し，表の列に制限はない~\cite[Definition~37]{Gallot2021}．その\emph{指数が高々$k$}とは，公理に表の列を適用して得られる全ての語が非終端記号の出現を高々$k$個しか含まないことをいう．

\begin{lemma}[ランク1のMCFLは有限指数EDT0Lである]
\label{lem:rank-one-edt0l}
$\mathsf{MCFL}(f',1)$の任意の言語は指数高々$f'$のEDT0L系をもつ．
\end{lemma}

\begin{proof}
$H$をファンアウト$\le f'$，ランク$\le1$の縮約済み文法とする．各非終端記号$B$と成分$j$について記号$[B,j]$を用い，公理を$[S,1]$とし，各規則$A\to\rho(B)$（または$A\to\rho$）について，$[A,j]$を$\rho$の第$j$成分（$\xi_1^k$を$[B,k]$に置き換えたもの）に写し，他の記号を固定する表を用いる．各規則の子は高々一つで線形なので，到達可能な各文形式は単一の非終端記号の記号を各高々一度ずつ含み，非終端記号の出現は高々$f'$個である．表の列はどれも許され，左辺が現れない規則の表は恒等写像として働く．

\emph{導出から表の列へ．}$H$の導出は規則の連鎖$R_1,\ldots,R_m$であり，$R_j$の子は$R_{j+1}$の左辺，$R_m$は子をもたない．$j$に関する帰納法により，$[S,1]$に$R_1,\ldots,R_j$の表を順に適用すると，$R_j$の子$B$を展開せずに残し，その第$k$成分を$[B,k]$と書いた（使われない成分は現れない）導出の産出が得られる．$j=m$のときこれは導出の終端語である．

\emph{表の列から導出へ．}表の列が終端語$w$を生むとする．左辺が現在の文形式に現れない段階は恒等写像として働くので除く．残る段階$R_1,\ldots,R_j$では，各左辺は現れている記号の非終端記号，すなわち直前の規則の子（または$S$）なので，これらは$S$からの連鎖をなす．$R_j$が子をもたなければ，この連鎖は産出$w$の導出である．そうでなければ$R_j$の子$B$の記号$[B,k]$は残っていないので，$B$の成分は全て消去されている．$H$は縮約済みなので$B$は何らかのタプルを導出し，連鎖を$B$の任意の導出で延長すると，産出がやはり$w$である$H$の導出が得られる．
\end{proof}

以下で用いるのはEDT0Lへの所属だけである．指数の上界は，非分岐MCFLと有限指数EDT0L言語を同一視するGallotの結果~\cite[Theorem~29]{Gallot2021}と対応させるために記した．

\begin{theorem}[文脈自由でない本質的ランク2の型付き標的]
\label{thm:intrinsic-ranktwo-witness}
$L_{\mathrm{intr}}\in\mathsf{MCFL}(2,2)\setminus\mathrm{CFL}$であり，任意の$f\ge2$について$L_{\mathrm{intr}}\in\Cmcf{f}{s_{1,1}}\setminus\YSub{f}$であるが，任意の有限$f'$について$L_{\mathrm{intr}}\notin\mathsf{MCFL}(f',1)$である．よってその最小ファンアウトは2，有限ファンアウトの全ての表示にわたる最小規則ランクは2である．定理~\ref{thm:boundary-polydata}はこれに適用できるが，定理~\ref{thm:rankone-polydata}はどの有限ファンアウトでも適用できない．
\end{theorem}

\begin{proof}
\emph{文法．}ファンアウト1の規則$B\to\mathtt aB\mathtt bB\mid\mathtt bB\mathtt aB\mid\mathtt aB\mathtt b\mid\mathtt{ab}B\mid\mathtt{ab}\mid\mathtt bB\mathtt a\mid\mathtt{ba}B\mid\mathtt{ba}$は$L_B=\mathrm{MIX}_2^+$を与える．生成される語は全て釣り合っており，逆に釣り合った語は$|\cdot|_{\mathtt a}-|\cdot|_{\mathtt b}$が最初に0に戻る位置で$\mathtt ay\mathtt bz$または$\mathtt by\mathtt az$（$y,z$は釣り合っており空でもよい）と分解されるからである．$A\to(\mathtt u,\mathtt w)\mid(\mathtt u\xi_1^1\mathtt v,\mathtt w\xi_1^2)(A)$により$L_A=\{(\mathtt u^n\mathtt v^{n-1},\mathtt w^n)\}$となり，規則$S\to(\alpha\mathtt\#\xi_1^1\mathtt\$\xi_2^1\mathtt v\xi_2^2)(B,A)$（$\alpha\in P_{\mathrm{intr}}$）が$L_{\mathrm{intr}}$を生成する．文法は縮約済み良形（$A$の成分は$\mathtt v$で隔てられる）で，ファンアウトとランクは2である．

\emph{文脈自由でない．}$\mathtt{lpc\#ab\$}\mathtt u^+\mathtt v^+\mathtt w^+$と交わり接頭部を消すと$T_3^+$が得られる．

\emph{型付きの所属．}$L_{\mathrm{intr}}=Q_{\mathrm{intr}}\cap\Phi^{-1}(0)$である．ここで$Q_{\mathrm{intr}}:=P_{\mathrm{intr}}\mathtt\#\{\mathtt a,\mathtt b\}^+\mathtt\$\,\mathtt u^+\mathtt v^+\mathtt w^+$は狭義2-局所であり，対は$\vdash\mathtt l$，$\vdash\mathtt r$，$\mathtt{lp},\mathtt{lq},\mathtt{rp},\mathtt{pc},\mathtt{pd},\mathtt{qc},\mathtt{c\#},\mathtt{d\#}$，$\mathtt{\#}x$，$xy$，$x\mathtt\$$（$x,y\in\{\mathtt a,\mathtt b\}$），$\mathtt{\$u},\mathtt{uu},\mathtt{uv},\mathtt{vv},\mathtt{vw},\mathtt{ww},\mathtt w\dashv$である（最初の六つの文字対はちょうど$P_{\mathrm{intr}}$の語を許す）．$\Phi:\Sigma_{\mathrm{intr}}^*\to\mathbb Z^3$は$\mathtt a,\mathtt b,\mathtt u,\mathtt v,\mathtt w$を$(1,0,0),(-1,0,0),(0,1,0),(0,-1,1),(0,0,-1)$に，他の文字を$0$に写す．補題~\ref{lem:sl2-boundary-control}と命題~\ref{prop:abelian-coset-filter}により$f\ge2$で$L_{\mathrm{intr}}\in\Cmcf{f}{s_{1,1}}$である．

\emph{$\YSub{f}$の外．}$x=\mathtt p$，$y=\mathtt q$，$E=\mathtt l\blank_1\mathtt{c\#ab\$uvw}$，$F=\mathtt r\blank_1\mathtt{d\#ab\$uvw}$について，$E[x],E[y],F[x]\in L_{\mathrm{intr}}$だが$F[y]\notin L_{\mathrm{intr}}$である．境界型付けは$x$と$y$を区別する．

\emph{ランク1は不可能．}$\mathsf{MCFL}(f',1)$は文字から語への準同型で閉じている（テンプレートに適用すればよい）．$L_{\mathrm{intr}}\in\mathsf{MCFL}(f',1)$なら，$\mathtt a,\mathtt b$以外の文字を消すと$\mathrm{MIX}_2^+\in\mathsf{MCFL}(f',1)$となり，補題~\ref{lem:rank-one-edt0l}によりEDT0Lとなり，さらに（EDT0Lは和で閉じているので）$\mathrm{MIX}_2=\mathrm{MIX}_2^+\cup\{\eps\}$もEDT0Lとなる．しかし$\mathrm{MIX}_2$は$\mathbb Z$の語の問題であり，EDT0Lではない~\cite{BishopEtAl2026EDT0L}（\cite[Corollary~32]{Gallot2021}も参照）．補題~\ref{lem:rank-one-edt0l}はGallotによる非分岐MCFGの特徴づけ~\cite[Definition~38 and Theorems~28--29]{Gallot2021}と整合する．同文献では定義と定理の番号は別々に付けられている．
\end{proof}

二つの下界は別々の因子から生じ，非分岐性による下界は先行研究である．本稿に特有なのは，狭義局所的な接頭部と$(1,1)$境界型付けが標的を型付きクラスの中に置く一方で，代入含意がそれをYoshinakaの基準線の外に置く点である．




% ---------------------------------------------------------------------
% [P2] 4.6節 制限のない標準入力MCFG
% ---------------------------------------------------------------------
\subsection{制限のない標準入力MCFG}
\label{subsec:unrestricted-source}

入力文法は良形である必要はなく，消去・置換・融合規則や空成分を許す．そのとき重要なのは，各非空成分パターンを実現するコストである．

\begin{definition}[成分パターン厚さ]
\label{def:component-pattern-thickness}
標準MCFG $H$の生産的な非終端記号$A$について，$\operatorname{supp}(\tuple u):=\{i:u_i\ne\eps\}$，$\mathsf{Pat}_H(A):=\{\operatorname{supp}(\tuple u):\tuple u\in L_A(H)\}$とし，
\[
 t_{\mathrm{pat}}(H):=
 \max_{\substack{A\text{ 生産的}\\I\in\mathsf{Pat}_H(A)}}
 \min_{\substack{\tuple u\in L_A(H)\\ \operatorname{supp}(\tuple u)=I}}
 \|\tuple u\|_1 ,
 \qquad
 \bar t_{\mathrm{pat}}(H):=\max\{1,t_{\mathrm{pat}}(H)\}
\]
とおく（$\max\emptyset:=0$）．$H$が縮約済みかつ$\lambda$-自由なら全成分パターンだけが現れ，$t_{\mathrm{pat}}(H)=t_H$である．
\end{definition}

\begin{proposition}[パターン厚さで制御される良形正規化]
\label{prop:pattern-thickness-transfer}
$H$を$L(H)\cap\Sigma^+\ne\emptyset$なるファンアウト$\le f$，ランク$\le r$の標準MCFG，$H_{\mathrm{red}}$をその刈り込み形とする．このとき同じファンアウト・ランク上界をもつ$L(H)\cap\Sigma^+$の縮約済み良形MCFG $G$で
\[
 \|G\|=O_{f,r}(\|H\|),\qquad
 t_G\le t_{\mathrm{pat}}(H_{\mathrm{red}}),\qquad
 \bar\theta_G=O_{f,r}\bigl(\|H\|\,\bar t_{\mathrm{pat}}(H_{\mathrm{red}})\bigr)
\]
なるものが存在する．
\end{proposition}

\begin{proof}
命題~\ref{prop:good-normalization}の$G$をとる．補題~\ref{lem:normalization-stages}により，$G$の各非終端記号の言語は，$H_{\mathrm{red}}$の台パターンのファイバーの和集合の，成分の削除・置換・連結による像である．これらは総長を増やさないので$t_G\le t_{\mathrm{pat}}(H_{\mathrm{red}})$である（付録~\ref{app:pattern-thickness}）．補題~\ref{lem:thickness-equivalence}(iii)が$\bar\theta_G$の上界を与える．
\end{proof}

\begin{corollary}[任意の標準MCFGからの特徴データ総長]
\label{cor:unrestricted-standard-mcfg-data}
$H$を$L\in\Cmcf{f}{s}$に対するファンアウト$\le f$，ランク$\le r$の任意の有限標準MCFGとし，$n_H:=\max\{1,\|H\|\}$，$t:=\bar t_{\mathrm{pat}}(H_{\mathrm{red}})$とおく．作用素$\mathcal B^{\le r}_{f,s}$と$\mathcal B_{f,s}$は$|K|=O(n_H)$かつ
\[
 \|K\|_+=
 \begin{cases}
 O(n_H^3t),&\text{ある}k,\ell\text{と有限群への準同型}g\text{（自明でもよい）について}\widehat s\preceq(\widehat s_{k,\ell},g)\text{のとき},\\
 O(n_H^4t),&\mathsf S\in\mathbf{LG}\text{のとき}
 \end{cases}
\]
なる共通の特徴標本$K\subseteq L$をもつ．定数は$f,r,s,\Sigma$（および$k,\ell,g$）に依存し，$H$には依存しない．
\end{corollary}

\begin{proof}
$L\subseteq\{\eps\}$なら$K=L$とする．そうでなければ命題~\ref{prop:pattern-thickness-transfer}により$\|G\|=O(n_H)$，$\bar\theta_G=O(n_Ht)$なる$G$が得られ，定理~\ref{thm:window-group-polydata}または定理~\ref{thm:local-group-polydata}を適用する．
\end{proof}


% ---------------------------------------------------------------------
% [P3] 付録C（一般MCFG表示に対するパターン厚さの移行）
% ---------------------------------------------------------------------
\section{一般MCFG表示に対するパターン厚さの移行}
\label{app:pattern-thickness}

\begin{proof}[命題~\ref{prop:pattern-thickness-transfer}における$t_G\le t_{\mathrm{pat}}(H_{\mathrm{red}})$の証明]
$H_{\mathrm{red}}$の非終端記号$A$と実現されるパターン$I\in\mathsf{Pat}_{H_{\mathrm{red}}}(A)$について$F_{A,I}:=\{\tuple u\in L_A(H_{\mathrm{red}}):\operatorname{supp}(\tuple u)=I\}$とおく．これは長さ高々$t_{\mathrm{pat}}(H_{\mathrm{red}})$のタプルを含み，$L_A(H_{\mathrm{red}})=\bigcup_IF_{A,I}$である．射影$\pi_J$，置換$\sigma$，融合$\mu$の合成を\emph{許容写像}と呼ぶ．これは総長を増やさず，$\phi(\tuple u)$の台は$\tuple u$の台だけで決まる．$H_{\mathrm{red}}$に補題~\ref{lem:normalization-stages}の段階を適用する．(a)と(b)の後，各非終端記号$X$は元の$A$と許容写像$\phi_X$について$L_X=\phi_X(L_A)$を満たす．段階(c)は実現される$I'$について$L_{X^{I'}}=\pi_{I'}\{\tuple u\in L_X:\operatorname{supp}(\tuple u)=I'\}$なる$X^{I'}$を作る．像の台は台だけで決まるので，これはある実現される$I$について$\pi_{I'}\phi_X(F_{A,I})$を含む．段階(d)は言語を融合による像に置き換え，最後の刈り込みは残る言語を変えない．よって$G$の各非終端記号$Y$は，ある$A$，実現される$I$，許容写像$\psi$について$\psi(F_{A,I})\subseteq L_Y$を満たし，$\min_{\tuple u\in L_Y}\|\tuple u\|_1\le t_{\mathrm{pat}}(H_{\mathrm{red}})$である．
\end{proof}

\begin{proposition}[$\lambda$-自由な入力の良形変換における厚さ]
\label{prop:structural-good-thickness}
$G_{\mathrm{in}}$を，消去規則を含んでもよいファンアウト$\le f$，ランク$\le r_{\max}$の縮約済み$\lambda$-自由MCFGで，$L(G_{\mathrm{in}})\cap\Sigma^+\ne\emptyset$とする．良形正規化は同じ上界をもつ縮約済み良形な$G$で
\[
  \|G\|=O_{f,r_{\max}}(\|G_{\mathrm{in}}\|),\qquad
  t_G\le t_{G_{\mathrm{in}}},\qquad
  \bar\theta_G=O_{f,r_{\max}}\bigl(\|G_{\mathrm{in}}\|\max\{1,t_{G_{\mathrm{in}}}\}\bigr)
\]
なるものを与える．
\end{proposition}

\begin{proof}
縮約済み$\lambda$-自由な$G_{\mathrm{in}}$では全成分パターンだけが現れるので$t_{\mathrm{pat}}(G_{\mathrm{in}})=t_{G_{\mathrm{in}}}$である．命題~\ref{prop:pattern-thickness-transfer}を適用する．
\end{proof}

したがって消去規則だけでは厚さの移行は妨げられない．妨げとなり得るのは真の空成分である（命題~\ref{prop:empty-component-blowup}）．


% ---------------------------------------------------------------------
% [P4] 付録F：prop:empty-component-blowup と後続段落
% ---------------------------------------------------------------------
\begin{proposition}[空成分の除去は厚さを爆発させ得る]
\label{prop:empty-component-blowup}
各$n\ge1$について，$t_{E_n}\le2$，$\theta_{E_n}=O(n)$であるサイズ$O(n)$の縮約済みファンアウト2，ランク2のMCFG $E_n$で，その段階(c)の変換が厚さ$2^n$以上の非終端記号と規則を含むものが存在する．特に$t_{\mathrm{pat}}(E_n)\ge2^n$である．
\end{proposition}

\begin{proof}
$E_n$の規則を$A_0\to(\mathtt a,\eps)$，$1\le i\le n$について$A_i\to(\eps,\mathtt c)\mid(\xi_1^1\xi_2^1,\xi_1^2\xi_2^2)(A_{i-1},A_{i-1})$，$S\to(\xi_1^1\mathtt d\,\xi_1^2)(A_n)$とする．これは縮約済みで，非消去，非置換，非融合である．各$A_i$を$(\eps,\mathtt c)$で（$A_0$は$(\mathtt a,\eps)$で）完成させれば$t_{E_n}\le2$であり，長さ高々$n$の経路に沿って$\theta_{E_n}=O(n)$である．台$\{1\}$の$A_i$-タプルは$(\eps,\mathtt c)$を使えないので$(\mathtt a^{2^i},\eps)$である．よって$A_n^{\{1\}}$と規則$S\to(\xi_1^1\mathtt d)(A_n^{\{1\}})$の厚さは$2^n$以上である．
\end{proof}

これは命題~\ref{prop:typed-thickness-exponential}とは異なる．ここでは元の表示の台パターンが指数的に長いタプルを強制し，あちらでは型付き精密化がそうする．どちらも特徴データの下界ではない．



% ---------------------------------------------------------------------
% [P5] 4.5節末：精密化厚さの境界の段落，cor:refinement-dichotomy，rem:refinement-vs-data
% ---------------------------------------------------------------------
\paragraph{完全精密化の厚さに関する代数的境界．}
$f\ge1$について，擬多様体$\mathbf V$が\emph{ファンアウト$f$で多項式精密化厚さをもつ}とは，任意の空でない有限$\Sigma$，任意の$r_{\max}\ge1$，$\mathsf S\in\mathbf V$なる任意の全射$s:\Sigma^+\twoheadrightarrow\mathsf S$について，多項式$p=p_{f,r_{\max},s,\Sigma}$が存在し，$\Sigma$上のファンアウト$\le f$，ランク$\le r_{\max}$の全ての縮約済み良形文法$G$が$\theta_{\mathrm{typ}}(G,s)\le p(\|G\|,\bar\theta_G)$を満たすことをいう．

\begin{corollary}[精密化厚さの二分法]
\label{cor:refinement-dichotomy}
任意の$f\ge1$と有限半群の任意の擬多様体$\mathbf V$について，$\mathbf V$がファンアウト$f$で多項式精密化厚さをもつのは$\mathbf V\subseteq\mathbf{LG}$のとき，かつそのときに限る．$\mathbf V\subseteq\mathbf{LG}$なら$\theta_{\mathrm{typ}}(G,s)=O_{f,r_{\max},s,\Sigma}(\|G\|^2\bar\theta_G)$ととれる．$G$を$L\in\Cmcf{f}{s}$なる$L\cap\Sigma^+$の表示に限っても同値性は保たれる．また$\theta_{\mathrm{typ}}$を一貫して$t_{\mathrm{typ}}$に置き換えても同じ同値性が成り立つ．$\mathsf S\notin\mathbf{LG}$なる個々の型付けについて，多項式上界はファンアウト1・規則ランク2ですでに成り立たない．
\end{corollary}

\begin{proof}
$\mathsf S\in\mathbf{LG}$なら，定理~\ref{thm:local-group-polydata}の証明が$\beta_\lambda=O(\|G\|^2\bar\theta_G)$で$(\mathrm R_\beta^\lambda)$を与え，代入可能性を用いない補題~\ref{lem:witness-length-template}(i),(iii)により$\theta_{\widetilde G_0}\le2\bar\theta_G+r_{\max}\beta_\lambda$である．

$\mathsf S\notin\mathbf{LG}$なら，定理~\ref{thm:lg-characterization}(ii)により部分半群$\{e,z\}\cong U_1$，すなわち$ez=ze=z$なるべき等元$e\ne z$が存在する．全射性から$s(u)=e$，$s(v)=z$なる$u,v\in\Sigma^+$を選ぶ．命題~\ref{prop:typed-thickness-exponential}(a)により，$\|G_n\|=O(n)$，$\bar\theta_{G_n}=O(1)$，$\theta_{\mathrm{typ}}(G_n,s)\ge2^n$なるファンアウト1，ランク2の文法$G_n$が得られ，これらは$\Sigma^+=\widehat s^{-1}(\mathsf S)\in\Cmcf{f}{s}$を表示する．最後に，$\mathbf V\not\subseteq\mathbf{LG}$なら$\mathsf S\in\mathbf V\setminus\mathbf{LG}$を選べばよい．これはアルファベット$\Sigma:=\mathsf S$上の評価型付け$\mathsf S^+\to\mathsf S$の正語像である．$t_{\mathrm{typ}}$について，肯定方向は補題~\ref{lem:witness-length-template}(i)と同じ$(\mathrm R_\beta^\lambda)$評価から，否定方向は命題~\ref{prop:typed-thickness-exponential}(a)から従う．
\end{proof}

\begin{remark}[精密化の障害が現れる場所]
\label{rem:refinement-vs-data}
命題~\ref{prop:typed-thickness-exponential}の文法はいずれも$\Sigma^+$を生成する．固定した正語像$\mathsf S$の各元には長さ高々$|\mathsf S|$の正語代表がある．というのも，最短代表の空でない接頭積が反復すれば語を短縮できるからである．よって
$\mu^\star(\Sigma^+,s)\le|\mathsf S|$だが，$t_{\mathrm{typ}}(G_n,s)\ge2^n$である．指数的な厚さは開始記号ではなく，精密化の内部に存在する．したがってこれだけでは特徴データ長の下界にはならず，そのためには異なる標的言語を区別できないことを示す必要がある．また，精密化が厚くても短い分布的代表は存在し得る（命題~\ref{prop:mcfg-relaxed-witnesses}）．
\end{remark}


% ---------------------------------------------------------------------
% [P6] 4.4節：thm:boundary-polydata
% ---------------------------------------------------------------------
\begin{theorem}[接頭辞--接尾辞決定的型付けの多項式データ]
\label{thm:boundary-polydata}
$s\preceq s_{k,\ell}$，$L\in\Cmcf{f}{s}$，$L(G)=L\cap\Sigma^+$とする．このとき$\Wit(\widetilde G_0)\cup(\{\eps\}\cap L)$は$\mathcal B^{\le r_{\max}}_{f,s}$と$\mathcal B_{f,s}$の両方の特徴標本であり，$|\Wit(\widetilde G_0)|\le C_s(G)=O(\|G\|)$かつ
\[
 \|\Wit(\widetilde G_0)\|_+\le K_{k,\ell}\,\|G\|^2\bar\theta_G,\qquad
 K_{k,\ell}:=|\mathsf S^\sharp|^{fr_{\max}}
   \bigl(3+r_{\max}(r_{\max}+1)f(k+\ell+1)|M_{k,\ell}|^f\bigr)
\]
である．
\end{theorem}

\begin{proof}
系~\ref{lem:boundary-short-canonical}と，$\lambda=\widehat s_{k,\ell}$，$b=b_{k,\ell}$とした補題~\ref{lem:witness-length-template}(i),(ii),(iv)から，特徴標本であることと$\|\Wit\|_+\le C_s(G)(1+2\bar\theta_G+r_{\max}b_{k,\ell})$が従う．$C_s(G)\le|\mathsf S^\sharp|^{fr_{\max}}\|G\|$，$1+2\bar\theta_G\le3\|G\|\bar\theta_G$，$|V|\le\|G\|$を代入する．
\end{proof}



% ---------------------------------------------------------------------
% [P7] 4.4節：thm:window-abelian-polydata と直後の段落
% ---------------------------------------------------------------------
\begin{theorem}[窓・アーベル群型付けの多項式データ]
\label{thm:window-abelian-polydata}
ここで$\exp({\mathbb G})$は有限群${\mathbb G}$の指数（全要素の位数の最小公倍数）を表す．定理~\ref{thm:window-group-polydata}で${\mathbb G}$が可換なら$c_{{\mathbb G},f}=\mathsf d({\mathbb G}^f)=O(\exp({\mathbb G})\,f\log|{\mathbb G}|)$（備考~\ref{rem:davenport-size}）であり，$K_{k,\ell,{\mathbb G}}$は因子$|\mathsf S^\sharp|^{fr_{\max}}$と$(|M_{k,\ell}||{\mathbb G}|)^f$を除いて$f$の多項式である．非可換な${\mathbb G}$では$c_{{\mathbb G},f}=M_{{\mathbb G},f}-1$は$f$について二重指数的に増大する．
\end{theorem}

\begin{proof}
補題~\ref{lem:block-replacement}(a)による．
\end{proof}

ファンアウト1では，固定窓に対する多項式時間・thick-data学習はYoshinaka~\cite[Theorem~1]{Yoshinaka2008}による．上の上界は縮約済み良形表示に相対的なものであり，正規化後にその結果を改良するとは主張しない．新しいのは，窓の値と順序付き群値を，交錯する高々$f$個の成分にわたって\emph{同時に}保つ点である．この高ランクの領域は言語のレベルで必要である．定理~\ref{thm:intrinsic-ranktwo-witness}は，最小ファンアウト2かつ最小規則ランク2の文脈自由でない$(1,1)$-型付き標的を与え，ランク1の場合は定理~\ref{thm:rankone-polydata}が扱う．

% ---------------------------------------------------------------------
% [P8] 4.5節：cor:mcfg-li-group-polydata
% ---------------------------------------------------------------------
\begin{corollary}[多項式データの代数的十分条件]
\label{cor:mcfg-li-group-polydata}
$\mathsf S\in\mathbf{LI}\vee\mathbf{Grp}$なら，任意の$L\in\Cmcf{f}{s}$と$L(G)=L\cap\Sigma^+$を満たす上述の任意の$G$について，$\Wit(\widetilde G_0)\cup(\{\eps\}\cap L)$は$\mathcal B^{\le r_{\max}}_{f,s}$と$\mathcal B_{f,s}$の両方の特徴標本であり，要素数は$O(\|G\|)$，総長は$O_{f,r_{\max},s,\Sigma}(\|G\|^2\bar\theta_G)$である．
\end{corollary}

\begin{proof}
補題~\ref{lem:mcfg-li-group-window}と定理~\ref{thm:window-group-polydata}による．$k$，${\mathbb G}$，$g$は$s$のみに依存する．
\end{proof}

% ---------------------------------------------------------------------
% [P9] 4.2節：一般形(a)を含む元の prop:typed-thickness-exponential
% ---------------------------------------------------------------------
\begin{proposition}[完全生成語型精密化の指数的増大]
\label{prop:typed-thickness-exponential}
\begin{enumerate}[label=\textup{(\alph*)},leftmargin=*]
\item $e\ne z$を$ez=ze=z$なる$\mathsf S$のべき等元とし，$u,v\in\Sigma^+$を$s(u)=e$，$s(v)=z$とする．$n\ge1$について$G_n$をファンアウト1，ランク2の文法
\[
 S\to SS\mid a\ (a\in\Sigma)\mid A_0\mid\cdots\mid A_n,\qquad
 A_0\to u,\qquad
 A_i\to A_{i-1}A_{i-1}\mid v\quad(1\le i\le n)
\]
とする（$u,v$は定数テンプレート）．このとき$G_n$は縮約済み良形，$L(G_n)=\Sigma^+$，$\|G_n\|=O_{\Sigma,u,v}(n)$，$\theta_{G_n}\le\max\{2,2|u|,2|v|\}$であり，$t_{\mathrm{typ}}(G_n,s)\ge|u|2^n$かつ$\theta_{\widetilde G_{n,0}}\ge|u|2^n$である．
\item 特に，$\Sigma=\{\mathtt a,\mathtt c\}$，$s_{\mathtt c}(\mathtt a)=1$，$s_{\mathtt c}(\mathtt c)=\zeta$で定まる$\{1,\zeta\}$（$\zeta^2=\zeta$）への二元型付け$s_{\mathtt c}$について，$u=\mathtt a$，$v=\mathtt c$とすれば$\|G_n\|=O(n)$，$\theta_{G_n}\le2$，$\theta_{\widetilde G_{n,0}}\ge2^n$である．
\end{enumerate}
\end{proposition}

\begin{proof}
全てのテンプレートは非空で各変数を順に一度ずつ使うので$G_n$は良形である．縮約済みであり，$S\to SS\mid a$により$\Sigma^+$を生成する．各$A_i$を$v$（$i\ge1$）または$u$（$i=0$）で完成させれば，各規則は長さ高々$\max\{2,2|u|,2|v|\}$の成功した使用をもつ．$A_n$の生成語は葉の値$s(u)=e$，$s(v)=z$の積であり，$e^2=e$かつ$z$は$\{e,z\}$で吸収的なので，その型が$e$であるのは$v$を使う葉がないときに限る．よって$(A_n)_{(e)}$の生成語は$u^{2^n}$だけであり，$t_{\mathrm{typ}}(G_n,s)\ge|u|2^n$が従う．型付き規則$S_{(e)}\to(A_n)_{(e)}$は（$S\to A_n$を介して）残り，これを使う成功導出の長さは$|u|2^n$以上である．(b)では$e=1$，$z=\zeta$である．
\end{proof}

% ---------------------------------------------------------------------
% [P10] 4.1節：厚さ尺度の表 tab:thickness-scales
% ---------------------------------------------------------------------
\begin{table}[htbp]
\caption{厚さ尺度：二つの入力尺度，中心となる型付き尺度と二つの補助量．}
\label{tab:thickness-scales}
\centering
\footnotesize
\setlength{\tabcolsep}{3.5pt}
\renewcommand{\arraystretch}{1.18}
\begin{tabularx}{\linewidth}{@{}p{0.155\linewidth}p{0.405\linewidth}X@{}}
\hline
尺度 & 測るもの／役割 & 比較または障害\\
\hline
$\theta_G$ & 各規則を使う最短受理文の長さの最大値．\emph{縮約済み良形$G$の主尺度}． & $t_G\le\theta_G\le 2\|G\|\max\{1,t_G\}$（補題~\ref{lem:thickness-equivalence}）．\\
$t_G$ & 非終端記号の最短生成語総長の最大値．良形での証明用尺度． & 縮約済み良形表示では$\theta_G$と多項式的に同値．\\
$t_{\mathrm{pat}}(H)$ & \emph{実現される非空成分パターン}ごとの最短生成語長の最大値．\emph{任意の$H$の主尺度}． & $t_H\le t_{\mathrm{pat}}(H)$，縮約済み$\lambda$-自由$H$では等号．族$E_n$は$\theta_{E_n}=O(n)$だが$t_{\mathrm{pat}}(E_n)\ge2^n$（命題~\ref{prop:empty-component-blowup}）．\\
$t_{\mathrm{typ}}$ & 開始記号以外の型付き非終端記号が導出する最短タプル総長の最大値．定理~\ref{thm:data-from-typed-thickness}の\emph{中心パラメータ}． & $t_{\mathrm{typ}}\le\theta_{\mathrm{typ}}\le2\bar\theta_G+r_{\max}t_{\mathrm{typ}}$（補題~\ref{lem:typed-thickness-equivalence}）．\\
$\theta_{\mathrm{typ}}$ & 刈り込んだ完全生成語型精密化$\widetilde G_0$の規則厚さ． & $\theta_{\mathrm{typ}}\ge\theta_G$．$\theta_G\le2$でも$\ge2^n$となり得る（命題~\ref{prop:typed-thickness-exponential}）．\\
$\kappa_G$ & 各型付き分布同値類の最短代表の長さの最大値．アーベル型フィルタ付き標的の補助量． & $\kappa_G\le\theta_{\mathrm{typ}}$．代表は型付き非終端記号から導出できるとは限らない（命題~\ref{prop:mcfg-relaxed-witnesses}）．\\
\hline
\end{tabularx}
\end{table}

% ---------------------------------------------------------------------
% [P11] 6.1節：統合前の prop:feature-cross-typed と cor:feature-relation-family
% ---------------------------------------------------------------------
\begin{proposition}[有限素性の認可を伴う交差依存]
\label{prop:feature-cross-typed}
$\Sigma_{\mathrm{feat}}=\{\mathtt p_0,\mathtt p_1,\mathtt q_0,\mathtt q_1,\mathtt a,\mathtt b,\mathtt c,\mathtt d\}$，$\mathcal R_{\mathrm{feat}}=\{(0,0),(0,1),(1,0)\}$，
\[
 L_{\mathrm{feat}}=\{\mathtt a^m\mathtt p_i\mathtt b^n\mathtt c^m\mathtt q_j\mathtt d^n:
 m,n\ge1,\ (i,j)\in\mathcal R_{\mathrm{feat}}\}
\]
とする．$g:\Sigma_{\mathrm{feat}}^*\to\mathbb Z_2^2$を$\mathtt p_1\mapsto(1,0)$，$\mathtt q_1\mapsto(0,1)$，他の文字$\mapsto0$で定め，$s:=(s_{1,1},g)|_{\Sigma_{\mathrm{feat}}^+}$（像への値域制限）とする．このとき$L_{\mathrm{feat}}\in\mathsf{MCFL}(2,2)\setminus\mathrm{CFL}$であり，任意の$f\ge2$について$L_{\mathrm{feat}}\in\Cmcf{f}{s}\setminus\YSub{f}$である．さらに$L_{\mathrm{feat}}$は$(1,s_{1,1})$-タプル代入可能でない．
\end{proposition}

\begin{proof}
\emph{文法．}ファンアウト2の$A,B$と規則$A\to(\mathtt a,\mathtt c)\mid(\mathtt a\xi_1^1,\mathtt c\xi_1^2)(A)$，$B\to(\mathtt b,\mathtt d)\mid(\mathtt b\xi_1^1,\mathtt d\xi_1^2)(B)$，$(i,j)\in\mathcal R_{\mathrm{feat}}$についての$S\to(\xi_1^1\mathtt p_i\xi_2^1\xi_1^2\mathtt q_j\xi_2^2)(A,B)$により，$L_A=\{(\mathtt a^m,\mathtt c^m)\}$，$L_B=\{(\mathtt b^n,\mathtt d^n)\}$，$L(G)=L_{\mathrm{feat}}$を得る．文法は良形で，ファンアウトとランクは2である．$\mathtt a^+\mathtt p_0\mathtt b^+\mathtt c^+\mathtt q_0\mathtt d^+$と交わり標識を消すと，古典的な文脈自由でない言語$\{\mathtt a^m\mathtt b^n\mathtt c^m\mathtt d^n\}$が得られる．

\emph{型付き代入可能性．}$Q=\mathtt a^+\{\mathtt p_0,\mathtt p_1\}\mathtt b^+\mathtt c^+\{\mathtt q_0,\mathtt q_1\}\mathtt d^+$（狭義2-局所），$K_{\mathrm{feat}}:=\{w\in\Sigma_{\mathrm{feat}}^+:g(w)\in\mathcal R_{\mathrm{feat}}\}$とし，$\varphi$を$\mathtt a,\mathtt c,\mathtt b,\mathtt d\mapsto(1,0),(-1,0),(0,1),(0,-1)\in\mathbb Z^2$，標識$\mapsto0$とする．$Q$上では$w$の標識が$\mathtt p_i,\mathtt q_j$なら$g(w)=(i,j)$なので，$L_{\mathrm{feat}}=Q\cap K_{\mathrm{feat}}\cap\varphi^{-1}(0)$である．三つの言語は，それぞれ補題~\ref{lem:sl2-boundary-control}，定義~\ref{def:class-cmcf}の後の(c)（$K_{\mathrm{feat}}$は$s$-同値類の和集合），命題~\ref{prop:abelian-coset-filter}(ii)により$(f,s)$-タプル代入可能であり，命題~\ref{prop:abelian-coset-filter}(i)が交差を与える．

\emph{$g$なしでの破綻．}$x=\mathtt{ap_0b}$，$y=\mathtt{ap_1b}$，$E=\blank_1\mathtt{cq_0d}$，$F=\blank_1\mathtt{cq_1d}$について，$E[x],E[y],F[x]\in L_{\mathrm{feat}}$だが，$(1,1)\notin\mathcal R_{\mathrm{feat}}$なので$F[y]\notin L_{\mathrm{feat}}$である．ここで$s_{1,1}(x)=s_{1,1}(y)$だが$g(x)\ne g(y)$である．
\end{proof}

\begin{corollary}[一様な有限素性の認可]
\label{cor:feature-relation-family}
\emph{同じ}型付け$s$について，任意の$R\subseteq\{0,1\}^2$に対し，言語$L_R:=\{\mathtt a^m\mathtt p_i\mathtt b^n\mathtt c^m\mathtt q_j\mathtt d^n:m,n\ge1,\ (i,j)\in R\}$は全ての$f\ge2$について$\Cmcf{f}{s}$に属し，$R\ne\emptyset$なら文脈自由でない．型付けは族全体に対して一度だけ固定する．$\mathbb Z_p\times\mathbb Z_q$への$g$を用いれば，任意の固定した素性目録$\{0,\ldots,p-1\}\times\{0,\ldots,q-1\}$についても同様である．
\end{corollary}

\begin{proof}
$R$で添字付けられた開始規則を残し，$K_{\mathrm{feat}}$の$\mathcal R_{\mathrm{feat}}$を$R$に置き換え（依然として$s$-同値類の和集合），証明を繰り返す．非文脈自由性には$(i,j)\in R$について$\mathtt a^+\mathtt p_i\mathtt b^+\mathtt c^+\mathtt q_j\mathtt d^+$を用いる．
\end{proof}

% ---------------------------------------------------------------------
% [P12] 6.2節：rem:necessity-scope
% ---------------------------------------------------------------------
\begin{remark}[型付けクラスの分離が意味する範囲]
\label{rem:necessity-scope}
最初の二つの分離は補題~\ref{lem:flag-separation}によりすでにアリティ1で起こり，群値一致の分離はアリティ2の文脈を用いる．標的の型付け$s$を，その狭義2-局所な構造言語を認識する有限準同型$h$で精密化したものを$s_{\mathrm{ref}}:=(s,h)$とする．すると各標的は，ある受理集合$\mathsf{Acc}'$とアーベル準同型$\varphi$により$\widehat s_{\mathrm{ref}}^{-1}(\mathsf{Acc}')\cap\varphi^{-1}(0)$と書けるので，定理~\ref{thm:mcfg-abelian-filter-data}は\emph{別の}作用素$\mathcal B^{\le2}_{f,s_{\mathrm{ref}}}$について，アルファベットに依存する次数$3|\Sigma|+2$の多項式データを与える．これに対し窓・群上界と局所群上界は，元の固定型付けのクラス全体に$|\Sigma|$によらない次数で一様に成り立つ．
\end{remark}

% ---------------------------------------------------------------------
% [P14] 結論（改稿前）
% ---------------------------------------------------------------------
\section{結論}
\label{sec:conclusion}

本稿は，型付きMCFGに対する正例学習の質的原理を，与えられた表示に相対的な特徴標本の\emph{符号化総長}の上界へと変換した．アルファベット，ファンアウト，ランク，型付けを固定すると，縮約済み良形文法$G$について，定理~\ref{thm:data-from-typed-thickness}は主なクラス一様の符号化総長上界を，計算可能な$t_{\mathrm{typ}}(G,s)$と$N_{\mathrm{typ}}=O(\|G\|)$へ還元する．具体的には次が成り立つ．
\begin{itemize}[leftmargin=*]
\item 窓型付けと窓・有限群型付けは，非可換群と交錯するタプル成分を含めて$O(\|G\|^2\bar\theta_G)$を与える（定理~\ref{thm:boundary-polydata}，\ref{thm:window-group-polydata}）．
\item 局所群型付けは，古典的なsliding-factor座標と所有者別ブロック短縮を組み合わせて$O(\|G\|^3\bar\theta_G)$を与える（定理~\ref{thm:local-group-polydata}）．
\item ランク1は任意の型付けの下で$O(\|G\|^3)$を与え（定理~\ref{thm:rankone-polydata}），任意の標準入力文法は成分パターン厚さを通じて扱える（系~\ref{cor:unrestricted-standard-mcfg-data}）．
\end{itemize}
精密化厚さの二分法（系~\ref{cor:refinement-dichotomy}）は，完全生成語型精密化の厚さが多項式となる正確な境界を，任意のファンアウトで$\mathbf{LG}$に置く．これは証拠構成の限界であって，データの限界ではない．特徴データ長そのものについては，稀な型の下界（命題~\ref{prop:rare-type}）によって$\mathbf{LG}$が正確な境界となるのはファンアウト1のみであり，それも擬多様体レベルである（定理~\ref{thm:fanout-one-data}）．ファンアウト2以上で本稿が証明するのは上界のみである．指数下界の言語族はコンパニオンCFG論文~\cite[Proposition~6.4]{kuriyamaCFG}で導入されたものであり，本稿で用いる性質は命題~\ref{prop:bracketed-family}と付録~\ref{app:bracketed-family}で証明する．ファンアウト2では同じ言語族が任意の固定有限型付けのもとで代入可能性を満たさない（定理~\ref{thm:doubling-no-lift}）．一方，\ref{subsec:baseline}節の素性認可の族は，文脈自由でない交差依存言語に適した固定有限バイアスの例を与える．

\begin{openproblem}[高ファンアウトでの特徴データ長]
\label{op:fanout-two}
$f\ge2$，$r_{\max}\ge2$を固定する．正語像$\mathsf S\notin\mathbf{LG}$なる型付け$s$と多項式$p$で，任意の$L\in\Cmcf{f}{s}$と，$L\cap\Sigma^+$のファンアウト$\le f$，ランク$\le r_{\max}$の任意の縮約済み良形表示$G$について，\emph{固定された}作用素$\mathcal B^{\le r_{\max}}_{f,s}$が総長高々$p(\|G\|,\bar\theta_G)$の特徴標本をもつものは存在するか．

反対方向の問いとして，正語像が$U_1$である全射型付け$s$を固定する．ある固定した$f\ge2$，$r_{\max}\ge2$と，縮約済み良形表示$G_n$（ファンアウト$\le f$，ランク$\le r_{\max}$）をもつ標的$L_n\in\Cmcf{f}{s}$で，$\|G_n\|\bar\theta_{G_n}=n^{O(1)}$でありながら，$L_n$に対する$\mathcal B^{\le r_{\max}}_{f,s}$の\emph{全ての}特徴標本の総長が$n^{\omega(1)}$となる族は存在するか．これは指定した学習作用素に関する下界であり，直ちに全ての集合駆動型学習器の下界とはならない．精密化厚さの二分法（系~\ref{cor:refinement-dichotomy}）は，このいずれの特徴データ長の問いにも答えない．命題~\ref{prop:rare-type}により，$\mu^\star(L_n,s)=n^{\omega(1)}$なる言語族を構成できれば下界には十分であり，その場合は固定型付けのクラス全体で特徴標本をもつ任意の集合駆動型作用素に同じ下界が従う．定理~\ref{thm:doubling-no-lift}は，命題~\ref{prop:bracketed-family}の括弧付き倍増言語族をその候補から排除する．
\end{openproblem}

より焦点を絞った未証明の問題として，任意の固定型付け$s$と規則ランク上限$r_{\max}$について，多項式$p_{s,r_{\max},\Sigma}$が存在し，$L\in\Cmcf{2}{s}$のファンアウト$\le2$，ランク$\le r_{\max}$の全ての縮約済み良形表示$G$について$\mu^\star(L,s)\le p(\|G\|,\bar\theta_G)$となるかを問うことができる．これは局所群型付けでは式~\eqref{eq:root-thickness}と定理~\ref{thm:local-group-polydata}により成立し，定理~\ref{thm:mcfg-abelian-filter-data}のアーベル型フィルタ付き標的でも成立する．ファンアウト1では類似の一般命題は成り立たない．冪の移行補題とこれらの特殊な場合だけでは問いは解決しない．

\paragraph{指数のタイトさ．}
多項式の指数の各因子の由来は明示できる．$O(\|G\|^2\bar\theta_G)$において，一つ目の$\|G\|$は証拠語の個数を数えるものであり（$N_{\mathrm{typ}}\le|\mathsf S^\sharp|^{fr_{\max}}\|G\|$，定理~\ref{thm:data-from-typed-thickness}），二つ目は$\lambda$-極小導出の根から葉への経路長を節点ラベルの個数$|V||\mathsf N|^f$で抑えることから生じる（補題~\ref{lem:heavy-node-count}）．局所群上界$O(\|G\|^3\bar\theta_G)$のさらなる因子は，重い部分木をまず長さ$O(\|G\|\bar\theta_G)$の窓代表と交換することによる（系~\ref{lem:boundary-short-canonical}，補題~\ref{lem:local-group-shortest}）．これはこの二段階の論法に由来する人工的な因子かもしれない．同様に，非可換群で定数$M_{\mathbb G,f}$が$f$について二重指数的になるのはRamsey型の均質化（補題~\ref{lem:order-homogenize}）による．アーベル群の場合のように$\mathrm{poly}(|\mathbb G|^f)$個の所有者で足りるかどうかは未解決である．一方，$\bar\theta_G$は取り除けない．\ref{subsec:rank-one}節の$\{\mathtt c^{2^n}\}$に対するサイズ$O(n)$の文法は$\bar\theta_G=2^n$をもち，全ての特徴標本は$\mathtt c^{2^n}$を含む．$\mathsf S\in\mathbf{LG}$の型付けについて，本稿の例が強制するのは$\Omega(\bar\theta_G)$までであり，残る差は窓・群型付けで因子$\|G\|^2$，局所群型付けで因子$\|G\|^3$である．ある固定型付け（自明なものでもよい）について，縮約済み良形表示$G_n$をもつ標的の族で，特徴標本の総長が$\Omega(\|G_n\|\,\bar\theta_{G_n})$を要するものは存在するか．稀な型の議論からはこれは得られない．固定型付けでは，式~\eqref{eq:root-thickness}によりその下界は高々$\sum_{m\in M_L}(\mu_m(L)+1)\le|\mathsf S|(t_{\mathrm{typ}}+1)$だからである．因子$\|G\|$を得るには，それぞれが二つの標的を区別する$\Theta(\|G\|)$個の語が必要であり，固定したアルファベットと型付けの下ではそれらの語を符号化せざるを得ない．すると同じ型の因子にわたる作用素の一般化により，族全体が少数の語から再構成されやすくなる．

したがって未解決なのは，固定した型付き標的クラスの全体に対する，特定の再構成作用素の特徴データの複雑さであり，個々の文法で完全生成語型精密化が厚くなるか否かではない．


